% BHU PYQ -- Professional Exam Template
\documentclass[11pt,a4paper]{article}

\usepackage[a4paper,top=2.5cm,bottom=2.5cm,left=2.8cm,right=2.8cm,headheight=14pt]{geometry}
\usepackage{amsmath,amssymb}
\usepackage[T1]{fontenc}
\usepackage[utf8]{inputenc}
\usepackage{lmodern}
\usepackage{microtype}
\usepackage{fancyhdr}
\usepackage{enumitem}
\usepackage{hyperref}
\usepackage{lastpage}
\usepackage{parskip}

\hypersetup{colorlinks=true,urlcolor=black,linkcolor=black,
  pdftitle={B.Sc. (Hons.) Botany Sem-III BOB-301 -- BHU PYQ}}

\pagestyle{fancy}
\fancyhf{}
\renewcommand{\headrulewidth}{0.4pt}
\renewcommand{\footrulewidth}{0.4pt}
\fancyhead[L]{\small\textit{Banaras Hindu University}}
\fancyhead[R]{\small\textit{Botany Paper: BOB-301}}
\fancyfoot[C]{\small Page \thepage\ of \pageref{LastPage}}

\newlist{parts}{enumerate}{2}
\setlist[parts,1]{label=(\alph*),leftmargin=*,itemsep=4pt,topsep=3pt}
\setlist[parts,2]{label=(\roman*),leftmargin=*,itemsep=2pt,topsep=2pt}
\newcommand{\pts}[1]{\hfill\textit{[#1]}}

\begin{document}

\begin{center}
  {\large\bfseries BANARAS HINDU UNIVERSITY}\\[2pt]
  {\normalsize B.Sc. (Hons.) Semester III Examination, 2016-17}\\[6pt]
  \rule{0.8\linewidth}{0.5pt}\\[6pt]
  {\large\bfseries Botany}\\[4pt]
  {\large\bfseries Paper: BOB-301 --- Plant Ecology \& Physiology}\\[4pt]
  {\normalsize Time: Three hours \hfill Max. Marks: 70}
\end{center}

\rule{\linewidth}{0.4pt}

\smallskip
\noindent\textit{Note: Answer five questions in all, selecting two each from Sections A and B. Question No. 1 is compulsory. The figures in the right-hand margin indicate marks.}

\bigskip

\medskip\noindent\textbf{Question 1.} Answer/Define the following in not more than 50 words each:\pts{$1 \times 10 = 10$}
\begin{parts}
  \item Synecology
  \item Net primary productivity
  \item Allogenic succession
  \item Food web
  \item Ecological pyramid
  \item Guttation
  \item Phototropism
  \item What is apical dominance?
  \item Which products of light reaction of photosynthesis are used in the Calvin cycle?
  \item $\text{H}^+$-ATPase
\end{parts}

\bigskip
\begin{center}
  \textbf{SECTION -- A}
\end{center}

\medskip\noindent\textbf{Question 2.} Write critical notes on any three of the following:\pts{$5 \times 3 = 15$}
\begin{parts}
  \item Atmosphere
  \item Soil profile
  \item Qualitative analytical characters of community
  \item Biotic components of ecosystem
\end{parts}

\medskip\noindent\textbf{Question 3.} Define population. Discuss in brief the various characteristics of a population of organisms.\pts{15}

\medskip\noindent\textbf{Question 4.} Write descriptive notes on any two of the following:\pts{$7.5 \times 2 = 15$}
\begin{parts}
  \item Ecotypes
  \item Phosphorus cycle
  \item Universal model of energy flow
\end{parts}

\medskip\noindent\textbf{Question 5.} Define ecological succession. Discuss in detail the course of succession in a hydrosere.\pts{15}

\bigskip
\begin{center}
  \textbf{SECTION -- B}
\end{center}

\medskip\noindent\textbf{Question 6.} Give an account of the $\text{C}_4$ and CAM pathways of $\text{CO}_2$ fixation in plants.\pts{15}

\medskip\noindent\textbf{Question 7.} Write descriptive notes on any two of the following:\pts{$7.5 \times 2 = 15$}
\begin{parts}
  \item Phytochromes
  \item Oxidative phosphorylation
  \item Photorespiration
\end{parts}

\medskip\noindent\textbf{Question 8.} Write short notes on any three of the following:\pts{$5 \times 3 = 15$}
\begin{parts}
  \item Translocation of sugar
  \item Assimilatory reduction of nitrate
  \item Ion transport
  \item ABA-induced stomatal closure
\end{parts}

\medskip\noindent\textbf{Question 9.} What are plant growth hormones? How do they differ from plant growth regulators? Give an account of the physiological role of auxins.\pts{15}

\vfill
\rule{\linewidth}{0.4pt}
\begin{center}\small
\href{https://bhu-exam.vercel.app/}{Click Here}\\[4pt]\href{https://me-aryan.vercel.app/}{Click Here}\\[4pt]\href{https://quashan-me.vercel.app/}{Click Here}
\end{center}

\end{document}
